\section{Matriz de Competências}

\vspace{0.5cm}

\textbf{FUNÇÃO 1} \\
Compreender os fundamentos científicos e tecnológicos da Inteligência Artificial para aplicação em soluções computacionais.

\begin{table}[htb]
\centering
\small
% reduce row height to avoid "Float too large for page" errors
\renewcommand{\arraystretch}{1.1}
\begin{tabular}{|p{6cm}|p{10cm}|}
\hline
\textbf{Subfunção} & \textbf{Padrões de Desempenho} \\
\hline
Interpretar os fundamentos da Inteligência Artificial. & Identificar problemas passíveis de solução por IA; diferenciar as principais áreas da IA; reconhecer limitações dos modelos inteligentes; interpretar conceitos fundamentais utilizados em Machine Learning. \\
\hline
Relacionar conceitos matemáticos utilizados em modelos inteligentes. & Aplicar princípios de álgebra linear, estatística e probabilidade; interpretar distribuições estatísticas; analisar métricas básicas utilizadas em modelos de IA. \\
\hline
Preparar ambientes computacionais para desenvolvimento de aplicações de IA. & Configurar ambientes Python; utilizar Linux e ferramentas de desenvolvimento; utilizar Git; preparar ambientes virtuais e containers. \\
\hline
Aplicar princípios éticos, legais e de governança relacionados ao desenvolvimento de IA. & Atuar conforme LGPD; aplicar princípios de IA Responsável; reconhecer vieses algorítmicos; respeitar normas de segurança e governança. \\
\hline
Desenvolver aplicações em Python para Inteligência Artificial. & Desenvolver algoritmos; utilizar estruturas de dados; manipular arquivos; modularizar aplicações; documentar código. \\
\hline
Preparar bases de dados para aplicações de IA. & Coletar dados; realizar limpeza; tratar valores ausentes; normalizar dados; realizar análise exploratória. \\
\hline
Desenvolver modelos básicos de Machine Learning. & Selecionar algoritmos; treinar modelos supervisionados e não supervisionados; validar modelos; interpretar métricas de desempenho. \\
\hline
Desenvolver modelos de Deep Learning. & Construir redes neurais; configurar treinamento; aplicar CNNs, RNNs e Transformers; otimizar hiperparâmetros; interpretar resultados. \\
\hline
Elaborar Engenharia de Prompt e Engenharia de Contexto. & Desenvolver prompts estruturados; aplicar técnicas Zero-shot, Few-shot e Chain-of-Thought; organizar contexto; avaliar desempenho de prompts; aplicar boas práticas de segurança em prompts. \\
\hline
Desenvolver aplicações utilizando Grandes Modelos de Linguagem. & Integrar LLMs; utilizar embeddings; desenvolver aplicações com RAG; estruturar saídas em JSON; integrar ferramentas através de Function Calling e Tool Calling. \\
\hline
\end{tabular}
\caption{Matriz de Competências - Função 1}
\label{tab:matriz_funcao1}
\end{table}

\vspace{1cm}

\newpage

\textbf{FUNÇÃO 2} \\
Desenvolver sistemas computacionais baseados em Inteligência Artificial, Análise Preditiva e Ecossistemas IoT, atendendo às normas de qualidade corporativas, robustez metodológica, governança de dados e arquitetura segura.

\begin{table}[htb]
\centering
\small
% reduce row height to avoid "Float too large for page" errors
\renewcommand{\arraystretch}{1.1}
\begin{tabular}{|p{6cm}|p{10cm}|}
\hline
\textbf{Subfunção} & \textbf{Padrões de Desempenho} \\
\hline
Desenvolver aplicações utilizando Inteligência Artificial Generativa. & Integrar modelos generativos; desenvolver chatbots inteligentes; implementar RAG; utilizar bancos vetoriais; validar respostas produzidas pelos modelos. \\
\hline
Projetar arquiteturas de agentes inteligentes. & Identificar arquiteturas de agentes; modelar agentes reativos, deliberativos e baseados em objetivos; definir fluxos de execução. \\
\hline
Desenvolver sistemas autônomos baseados em LLMs. & Implementar planejamento; utilizar memória; integrar raciocínio estruturado; desenvolver workflows inteligentes. \\
\hline
Integrar ferramentas utilizando Model Context Protocol (MCP). & Desenvolver MCP Servers; consumir MCP Clients; integrar Resources, Prompts e Tools; validar interoperabilidade entre aplicações. \\
\hline
Desenvolver aplicações utilizando LangGraph e frameworks de agentes. & Construir grafos de execução; implementar nós, estados e transições; desenvolver workflows complexos; integrar LangGraph, CrewAI, AutoGen e OpenAI Agents SDK. \\
\hline
Desenvolver sistemas multiagentes. & Implementar agentes especializados; coordenar comunicação entre agentes; aplicar arquiteturas hierárquicas e colaborativas; monitorar execução distribuída. \\
\hline
Desenvolver aplicações inteligentes utilizando Engenharia de Software. & Projetar arquiteturas; desenvolver APIs REST; integrar bancos de dados; aplicar padrões de projeto; documentar aplicações. \\
\hline
Desenvolver aplicações Full Stack integradas à Inteligência Artificial. & Desenvolver interfaces; integrar backend e frontend; consumir APIs inteligentes; implantar aplicações web. \\
\hline
Desenvolver soluções de Análise Preditiva. & Construir pipelines analíticos; selecionar modelos; validar previsões; interpretar indicadores; gerar relatórios técnicos. \\
\hline
Desenvolver aplicações de Visão Computacional. & Preparar imagens; desenvolver modelos CNN e YOLO; realizar classificação, segmentação e detecção de objetos; implantar modelos em produção. \\
\hline
Desenvolver soluções IoT e IIoT integradas à Inteligência Artificial. & Integrar sensores e gateways; utilizar MQTT e OPC UA; coletar telemetria; desenvolver aplicações Edge Computing; integrar IA aos dispositivos. \\
\hline
Desenvolver sistemas distribuídos para processamento inteligente. & Implementar pipelines utilizando Kafka, Spark e Flink; utilizar bancos de séries temporais; integrar Edge e Cloud Computing. \\
\hline
Implantar modelos em ambientes produtivos utilizando MLOps e LLMOps. & Automatizar pipelines; versionar modelos; monitorar desempenho; detectar deriva; implantar modelos em containers e Kubernetes. \\
\hline
Aplicar observabilidade e governança em soluções de IA. & Implementar monitoramento; coletar métricas; utilizar tracing; aplicar auditoria, segurança, conformidade e IA Responsável. \\
\hline
Desenvolver projetos integradores de Inteligência Artificial. & Planejar soluções; integrar tecnologias estudadas; validar requisitos; documentar projetos; apresentar soluções técnicas conforme padrões profissionais. \\
\hline
\end{tabular}
\caption{Matriz de Competências - Função 2}
\label{tab:matriz_funcao2}
\end{table}
